\documentclass{article}

% \usepackage{neurips_2026}
\usepackage[nonanonymous]{neurips_2026}

\usepackage[utf8]{inputenc}
\usepackage[T1]{fontenc}
\usepackage{hyperref}
\usepackage{url}
\usepackage{booktabs}
\usepackage{amsfonts}
\usepackage{nicefrac}
\usepackage{microtype}
\usepackage{xcolor}
\usepackage{graphicx}
\usepackage{amsmath}
\usepackage{multirow}
\usepackage{subcaption}
\usepackage{placeins}

\setcounter{topnumber}{4}
\setcounter{bottomnumber}{4}
\setcounter{totalnumber}{6}
\renewcommand{\topfraction}{0.9}
\renewcommand{\bottomfraction}{0.9}
\renewcommand{\textfraction}{0.05}
\renewcommand{\floatpagefraction}{0.7}

\title{Benchmarking Data Augmentation Methods for Cancer Subtype Classification from Bulk RNA-seq}

\author{%
  Ishaan Gadkari \\
  Department of Computer Science \\
  Virginia Tech \\
  \texttt{ishaangadkari@vt.edu}
  \And
  Damien Stone \\
  Department of Computer Science \\
  Virginia Tech \\
  \texttt{damienstone23@vt.edu}
}

\begin{document}

\maketitle

% ============================================================
\begin{abstract}
% ============================================================
Cancer subtype classification from bulk RNA-sequencing data is a clinically important problem
complicated by small labeled cohorts, severe class imbalance, and high-dimensional feature spaces.
Data augmentation offers a principled strategy to mitigate these challenges, yet no systematic
comparison of augmentation methods exists across multiple tumor types using the same experimental
protocol. We benchmark six augmentation families against a no-augmentation baseline on three
TCGA cohorts, covering SMOTE and its variants (BorderlineSMOTE, ADASYN, random oversampling),
the conditional variational autoencoder (cVAE), the Wasserstein GAN with gradient penalty
(WGAN-GP), and Mixup. The cohorts are
TCGA cohorts: breast cancer (BRCA, five PAM50 subtypes, $n=1{,}218$), colon adenocarcinoma
(COAD, four CMS subtypes, $n=329$), and prostate adenocarcinoma (PRAD, four molecular subtypes,
$n=550$). A multi-layer perceptron trained on log$_2$(RSEM$+$1) gene expression serves as the
downstream classifier, evaluated by macro-averaged AUROC, macro F1, and per-class F1 across
five random seeds. Our main findings are: (1) interpolation-based methods, particularly
BorderlineSMOTE and ADASYN, yield the largest and most consistent macro F1 gains on BRCA
($+0.03$ over baseline); (2) generative deep-learning methods (cVAE, WGAN-GP) provide no
reliable benefit and frequently harm performance, especially on minority classes; (3) augmentation
benefit is dataset-dependent: COAD is largely saturated and PRAD shows no macro F1 improvement
from any augmentation; and (4) synthetic sample quality, measured by Maximum Mean Discrepancy and
gene-gene correlation preservation, does not predict downstream classifier benefit, suggesting that
distribution-matching is an insufficient proxy for augmentation utility. These findings provide
actionable guidance for practitioners and motivate the development of augmentation-aware evaluation
metrics beyond distributional fidelity.
\end{abstract}

% ============================================================
\section{Introduction}
% ============================================================

The molecular subtyping of tumors from bulk RNA-sequencing has transformative clinical implications:
it guides treatment selection, predicts recurrence risk, and enables patient stratification in
clinical trials. Representative examples include the PAM50 subtypes of breast cancer (Luminal A,
Luminal B, HER2-enriched, Basal-like, Normal-like)~\cite{Parker2009}, the Consensus Molecular
Subtypes (CMS1--4) of colorectal cancer~\cite{Guinney2015}, and the ERG/ETV/SPOP fusion-based
subtypes of prostate cancer~\cite{Cancer2015}. These classification tasks are fundamentally
multi-class problems in extremely high-dimensional input spaces ($\sim$20{,}000 genes), with sample
sizes that are modest by deep learning standards and often severe class imbalance.

Machine learning classifiers trained on such data face several compounding difficulties. First, the
ratio of features to samples is unfavorable: even the largest TCGA cohorts offer at most a few
thousand labeled samples, whereas the feature dimension is an order of magnitude larger. Second,
minority subtypes, such as the Normal-like subtype in BRCA ($n{=}36$) or the FOXA1 subtype in PRAD
($n{=}9$), are so infrequent that a classifier can achieve high overall accuracy while completely
ignoring them. Third, the high intrinsic correlations among gene expression features mean that
naive oversampling may fail to respect the true data manifold.

Data augmentation is a natural response to these challenges. In computer vision and natural
language processing, augmentation routinely produces large gains~\cite{Shorten2019,Wei2019}.
For tabular and genomic data, the evidence is more mixed. SMOTE~\cite{Chawla2002} and its variants
are widely used for class-imbalanced tabular problems, while deep generative models, including variational
autoencoders~\cite{Kingma2014} and generative adversarial networks~\cite{Goodfellow2014}, have
attracted attention for their potential to capture complex data distributions. Mixup~\cite{Zhang2018}
offers a simple, training-time interpolation strategy that has shown benefits in several domains.
Yet no prior work systematically compares these families across multiple cancer types under a
controlled experimental protocol, leaving practitioners without clear guidance on which method to
use for genomic classification.

This paper addresses that gap. Our contributions are as follows:
\begin{itemize}
  \item We conduct the first large-scale benchmark of six data augmentation families for cancer
    subtype classification from bulk RNA-seq across three distinct TCGA cohorts (BRCA, COAD,
    PRAD), totaling over 20 augmentation configurations evaluated across five random seeds.
  \item We demonstrate that interpolation-based methods (BorderlineSMOTE, ADASYN) consistently
    outperform deep generative approaches (cVAE, WGAN-GP) in macro F1, while being
    substantially simpler and cheaper to train.
  \item We show that synthetic sample quality metrics (MMD, gene-gene correlation) do not
    reliably predict downstream classifier benefit, challenging the common assumption that
    better-looking synthetic data is more useful data.
  \item We identify the Normal-like subtype of BRCA as a persistent failure mode for all methods,
    and the FOXA1 subtype of PRAD as uncorrectable by any augmentation strategy tested.
\end{itemize}

% ============================================================
\section{Related Work}
% ============================================================

\paragraph{Oversampling for imbalanced classification.}
SMOTE~\cite{Chawla2002} generates synthetic minority-class examples by interpolating between a
sample and one of its $k$-nearest neighbors in feature space. BorderlineSMOTE~\cite{Han2005}
focuses synthesis on samples near the decision boundary, while ADASYN~\cite{He2008} adaptively
weights synthesis by local class difficulty. These methods were originally developed for binary
classification but have been extended to multi-class settings via one-vs-rest or one-vs-one
decompositions. In the biomedical domain, SMOTE-family methods have been applied to clinical
tabular data~\cite{Blagus2013} and, more recently, to genomic features~\cite{Lemasson2022},
though systematic comparisons across multiple tumor types are lacking.

\paragraph{Deep generative models for data augmentation.}
Variational autoencoders (VAEs)~\cite{Kingma2014} learn a low-dimensional latent representation
that enables conditional generation via a label-conditioned decoder. Conditional VAEs
(cVAEs)~\cite{Sohn2015} have been applied to gene expression synthesis~\cite{Marouf2020,Way2018},
with some success in augmenting rare cell types. GANs~\cite{Goodfellow2014} offer an adversarial
alternative; WGAN-GP~\cite{Gulrajani2017} stabilizes training via a gradient penalty and has
become a standard baseline for tabular GAN methods. scGAN~\cite{Marouf2020} and related models
apply GAN-based generation to single-cell RNA-seq but have not been systematically validated on
bulk RNA-seq subtype classification. A notable limitation of prior generative model evaluations
is that they primarily assess distributional fidelity (e.g., MMD, FID) rather than downstream
classifier performance, an issue our work directly addresses.

\paragraph{Mixup and feature-space augmentation.}
Mixup~\cite{Zhang2018} trains on convex combinations of training examples and their labels,
encouraging linear behavior between training points. In tabular domains, Mixup has shown modest
but consistent gains~\cite{Guo2019}. Its interaction with class imbalance is underexplored:
mixing a minority-class example with a majority-class example produces a fractional label, which
may or may not benefit minority-class recall.

\paragraph{Augmentation in cancer genomics.}
Existing studies applying augmentation to bulk RNA-seq cancer classification are typically
single-cancer and single-method~\cite{Yoon2020,Levy2020}, making cross-method and cross-cancer
comparisons impossible. Our work provides the first controlled benchmark across these dimensions,
enabling principled method selection for future studies.

% ============================================================
\section{Methodology}
% ============================================================

\subsection{Problem Formulation}

Let $\mathcal{D} = \{(\mathbf{x}_i, y_i)\}_{i=1}^{N}$ denote a labeled dataset where
$\mathbf{x}_i \in \mathbb{R}^G$ is a log$_2$(RSEM$+$1)-normalized bulk RNA-seq expression
vector over $G$ genes and $y_i \in \{1, \ldots, C\}$ is a cancer subtype label. The task is
multi-class classification. We evaluate whether augmenting the training set
$\mathcal{D}_\text{train}$ with synthetic examples improves classifier performance on a held-out
test set $\mathcal{D}_\text{test}$.

\subsection{Datasets and Preprocessing}

We use three TCGA cohorts downloaded from the UCSC Xena Hub (log$_2$(RSEM$+$1) format, HiSeqV2
pipeline). Clinical and subtype labels were obtained from the cBioPortal datahub GitHub LFS
mirror. Full preprocessing details are given below; deviations from original data sources are
documented in Appendix~\ref{app:data}.

\textbf{BRCA} ($n{=}1{,}218$ samples, $G{=}20{,}530$ genes). PAM50 subtype labels from the
TCGA Pan-Cancer Atlas: Luminal A (499), Luminal B (197), Basal-like (171), HER2-enriched (78),
Normal-like (36). Samples with missing subtype labels ($n{=}103$) are excluded. This yields a
heavily imbalanced five-class problem with a 14:1 majority-to-minority ratio.

\textbf{COAD} ($n{=}329$ samples after filtering, $G{=}20{,}530$ genes). Consensus Molecular
Subtype (CMS) labels from the CRC Subtyping Consortium~\cite{Guinney2015}, matched by TCGA
barcode: CMS2 (220), CMS4 (144), CMS1 (76), CMS3 (72). Samples labeled NOLBL ($n{=}61$) are
excluded per consortium guidelines.

\textbf{PRAD} ($n{=}333$ samples, $G{=}20{,}530$ genes). Molecular subtype labels from the
TCGA 2015 Cell paper~\cite{Cancer2015}. Following the specification, we collapse the seven
original subtypes to four: ERG-fusion (152), Other (86$+$3$+$4$=$merged rare classes), SPOP
mutation (37), and FOXA1 mutation (9). FOXA1 is severely underrepresented at 2.7\% of samples.

For all cohorts, features are standardized to zero mean and unit variance on the training set,
with the same transformation applied to validation and test sets.

\subsection{Classifier Architecture}

We use a multi-layer perceptron (MLP) as the downstream classifier, consisting of three fully
connected layers (1024$\to$512$\to$256) with ReLU activations, batch normalization, and dropout
($p{=}0.3$). The output layer uses softmax over $C$ classes. The model is trained with the Adam
optimizer, cross-entropy loss, and early stopping on the validation macro F1. This architecture
was selected for its ability to handle high-dimensional tabular inputs without explicit feature
selection and to serve as a consistent baseline across all augmentation methods.

\subsection{Augmentation Methods}

We benchmark six augmentation families, applied only to the training set:

\textbf{SMOTE} interpolates each minority-class sample with one of its $k$ nearest neighbors in
feature space: $\tilde{\mathbf{x}} = \mathbf{x}_i + \lambda(\mathbf{x}_{nn} - \mathbf{x}_i)$,
$\lambda \sim \text{Uniform}(0,1)$. We also evaluate tuned variants that optimize $k$ on the
validation set.

\textbf{BorderlineSMOTE} restricts synthesis to minority samples near the class boundary
(``danger zone''), focusing capacity on the hardest regions of feature space.

\textbf{ADASYN} assigns higher synthesis weights to minority samples that are more frequently
misclassified by a $k$-NN classifier, adaptively concentrating augmentation where it is most needed.

\textbf{Random oversampling} replicates minority-class samples with replacement until classes are
balanced, serving as a simple baseline for oversampling.

\textbf{Conditional VAE (cVAE)} learns an encoder $q_\phi(\mathbf{z}|\mathbf{x},y)$ and
decoder $p_\theta(\mathbf{x}|\mathbf{z},y)$ with a class-conditioned prior, optimizing the
evidence lower bound (ELBO):
\begin{equation}
  \mathcal{L}_\text{ELBO} = \mathbb{E}_{q_\phi}[\log p_\theta(\mathbf{x}|\mathbf{z},y)]
    - \text{KL}(q_\phi(\mathbf{z}|\mathbf{x},y) \| p(\mathbf{z}|y)).
\end{equation}
Synthetic samples are generated by sampling $\mathbf{z} \sim p(\mathbf{z}|y)$ and decoding.

\textbf{WGAN-GP} trains a generator $G$ and critic $D$ under the Wasserstein distance with
gradient penalty~\cite{Gulrajani2017}:
\begin{equation}
  \mathcal{L} = \mathbb{E}[D(G(\mathbf{z}))] - \mathbb{E}[D(\mathbf{x})]
    + \lambda \mathbb{E}[(\|\nabla_{\hat{\mathbf{x}}} D(\hat{\mathbf{x}})\|_2 - 1)^2],
\end{equation}
where $\hat{\mathbf{x}}$ interpolates between real and generated samples. Labels are concatenated
to both generator input and critic input for conditional generation.

\textbf{Mixup} generates training examples as convex combinations:
$(\tilde{\mathbf{x}}, \tilde{y}) = (\lambda \mathbf{x}_i + (1-\lambda)\mathbf{x}_j,
\lambda y_i + (1-\lambda)y_j)$, $\lambda \sim \text{Beta}(\alpha, \alpha)$.

For SMOTE-family methods, we also evaluate a \emph{sweep} variant that varies the number of
synthetic samples ($N_\text{synth} \in \{100, 500, 1000, 2000, 3000\}$) per minority class to
characterize sensitivity to augmentation volume. Additionally, we include Random Forest (RF)
and XGBoost (XGB) classifiers without augmentation as alternative baselines.

\subsection{Evaluation Protocol}

All experiments are repeated over five random seeds controlling train/validation/test splits
(70\%/15\%/15\% stratified), model initialization, and stochastic training. Primary metrics are
macro-averaged AUROC (one-vs-rest), macro F1, and per-class F1. Results are reported as mean
$\pm$ standard deviation over seeds. The no-augmentation MLP serves as the primary baseline.

% ============================================================
\section{Experiments}
% ============================================================

\subsection{Datasets}

The three cohorts span a range of difficulty: BRCA has the largest sample size but the most
severe class imbalance; COAD has the smallest cohort but a more balanced class distribution
(minimum class fraction 15\%); PRAD has a moderate sample size but an extreme minority class
(FOXA1 at 2.7\%). This variation allows us to assess whether augmentation benefit is
modulated by dataset properties.

\subsection{Baseline Methods}

We compare against three baselines: (1) the MLP without any augmentation (\textit{none}),
(2) a Random Forest without augmentation (\textit{rf\_none}), and (3) XGBoost without
augmentation (\textit{xgb\_none}). This design separates the effect of the augmentation
strategy from the effect of the classifier architecture.

\subsection{Evaluation Metrics}

We report macro-averaged AUROC as the primary ranking metric, as it is threshold-independent
and appropriate for imbalanced multi-class settings. Macro F1 captures the balance between
precision and recall across classes. Per-class F1 exposes which subtypes benefit or are
harmed by augmentation. Synthetic sample quality is characterized post hoc by MMD (a measure
of distributional distance between real and synthetic samples) and the Pearson correlation
between gene-gene covariance matrices of real and synthetic data.

% ============================================================
\section{Results}
% ============================================================

\subsection{AUROC}

Figure~\ref{fig:roc} shows macro-averaged one-vs-rest ROC curves for the six representative
augmentation conditions on each dataset (seed 42). All methods achieve high AUROC on COAD
($\geq$0.980), reflecting the relatively easier classification structure of CMS subtypes.
On BRCA, WGAN-GP achieves the highest AUROC (0.959), followed by Borderline-tuned (0.955)
and no augmentation (0.954); cVAE underperforms at 0.921. On PRAD, cVAE (0.975) and Mixup
(0.976) lead, while WGAN-GP (0.915) and SMOTE-tuned (0.939) lag behind the no-augmentation
baseline (0.965).

\begin{figure}[htbp]
  \centering
  \includegraphics[width=0.95\linewidth]{roc_grid.png}
  \caption{Macro-averaged one-vs-rest ROC curves for six augmentation conditions on BRCA, COAD,
    and PRAD (seed 42). All methods achieve high AUROC on COAD; the most differentiation
    appears on BRCA and PRAD.}
  \label{fig:roc}
\end{figure}
\FloatBarrier

\subsection{Macro F1 and Accuracy}

Table~\ref{tab:main} summarizes macro F1 for all methods across the three cohorts (mean over
five seeds).

\begin{table}[htbp]
  \centering
  \caption{Macro F1 (mean $\pm$ std over 5 seeds) for representative augmentation methods.
    Best result per column in \textbf{bold}. $\Delta$ is improvement over no-augmentation MLP baseline.}
  \label{tab:main}
  \small
  \begin{tabular}{lcccccc}
    \toprule
    & \multicolumn{2}{c}{\textbf{BRCA}} & \multicolumn{2}{c}{\textbf{COAD}} & \multicolumn{2}{c}{\textbf{PRAD}} \\
    \cmidrule(lr){2-3}\cmidrule(lr){4-5}\cmidrule(lr){6-7}
    \textbf{Method} & Macro F1 & $\Delta$ & Macro F1 & $\Delta$ & Macro F1 & $\Delta$ \\
    \midrule
    None (MLP)         & 0.763 & ---   & 0.862 & ---   & 0.652 & --- \\
    RF (no aug)        & 0.725 & $-$0.038 & 0.862 & 0.000 & 0.645 & $-$0.007 \\
    XGB (no aug)       & 0.677 & $-$0.086 & 0.843 & $-$0.019 & \textbf{0.728} & $+$0.076 \\
    \midrule
    SMOTE              & 0.775 & $+$0.012 & 0.860 & $-$0.002 & 0.621 & $-$0.031 \\
    SMOTE (tuned)      & 0.778 & $+$0.015 & 0.865 & $+$0.003 & 0.619 & $-$0.033 \\
    BorderlineSMOTE    & \textbf{0.796} & $+$0.033 & \textbf{0.865} & $+$0.003 & 0.598 & $-$0.054 \\
    ADASYN             & 0.790 & $+$0.027 & 0.861 & $-$0.001 & --- & --- \\
    Random oversample  & 0.774 & $+$0.011 & 0.859 & $-$0.003 & 0.622 & $-$0.030 \\
    cVAE               & 0.662 & $-$0.101 & 0.851 & $-$0.011 & 0.622 & $-$0.030 \\
    WGAN-GP            & 0.681 & $-$0.082 & \textbf{0.865} & $+$0.003 & 0.629 & $-$0.023 \\
    Mixup              & 0.766 & $+$0.003 & 0.862 & 0.000 & 0.620 & $-$0.032 \\
    \bottomrule
  \end{tabular}
\end{table}

On \textbf{BRCA}, interpolation methods provide consistent improvement. BorderlineSMOTE achieves
the highest macro F1 (0.796, $+$0.033 over baseline) and ADASYN follows closely (0.790,
$+$0.027). In contrast, cVAE drops sharply to 0.662 ($-$0.101), the lowest of all methods.
WGAN-GP also underperforms baseline (0.681). Mixup provides only a marginal gain (0.766).

On \textbf{COAD}, differences across methods are small (range: 0.843--0.865), suggesting the
dataset is largely saturated at this classifier capacity. BorderlineSMOTE and WGAN-GP jointly
achieve the highest macro F1 (0.865), while XGBoost without augmentation is notably weaker
(0.843), consistent with gradient-boosted trees being less effective on this cohort.

On \textbf{PRAD}, no augmentation strategy improves over the MLP baseline in macro F1. XGBoost
without augmentation (0.728) substantially outperforms all MLP-based configurations, suggesting
that classifier choice dominates augmentation choice for this dataset. Figure~\ref{fig:macro_f1}
shows the full macro F1 ranking across all configurations, and Figure~\ref{fig:accuracy} shows
the corresponding accuracy rankings.

\begin{figure}[htbp]
  \centering
  \includegraphics[width=0.95\linewidth]{summary_macro_f1.png}
  \caption{Macro F1 (mean $\pm$ std over seeds) for all augmentation configurations on BRCA
    (left), COAD (center), and PRAD (right), ranked by performance. The dashed line indicates
    the no-augmentation baseline.}
  \label{fig:macro_f1}
\end{figure}

\begin{figure}[htbp]
  \centering
  \includegraphics[width=0.95\linewidth]{summary_accuracy.png}
  \caption{Accuracy (mean $\pm$ std over seeds) for all augmentation configurations on BRCA
    (left), COAD (center), and PRAD (right), ranked by performance.}
  \label{fig:accuracy}
\end{figure}
\FloatBarrier

\subsection{Per-class F1}

Figure~\ref{fig:per_class} shows per-class F1 averaged over seeds for all methods on each cohort.

\begin{figure}[htbp]
  \centering
  \begin{subfigure}[b]{0.30\linewidth}
    \includegraphics[width=\linewidth, height=0.75\textheight, keepaspectratio]{per_class_f1_brca.png}
    \caption{BRCA}
  \end{subfigure}
  \begin{subfigure}[b]{0.30\linewidth}
    \includegraphics[width=\linewidth, height=0.75\textheight, keepaspectratio]{per_class_f1_coad.png}
    \caption{COAD}
  \end{subfigure}
  \begin{subfigure}[b]{0.30\linewidth}
    \includegraphics[width=\linewidth, height=0.75\textheight, keepaspectratio]{per_class_f1_prad.png}
    \caption{PRAD}
  \end{subfigure}
  \caption{Per-class F1 (mean over seeds) for all augmentation methods. Each row is a method;
    columns are subtypes. BRCA Normal-like and PRAD FOXA1 are the most problematic classes.}
  \label{fig:per_class}
\end{figure}
\FloatBarrier

Several patterns emerge. In \textbf{BRCA}, the Normal-like subtype ($n{=}36$) is the hardest:
all SMOTE-family methods achieve F1 of 0.54--0.66, while cVAE collapses to 0.07 and WGAN-GP
to 0.31. The Luminal B subtype also shows pronounced variation across methods (0.43--0.63),
reflecting confusion with Luminal A. Basal-like achieves near-perfect F1 ($\geq$0.94) across
all methods, consistent with its distinct transcriptomic profile.

In \textbf{PRAD}, FOXA1 ($n{=}9$) achieves F1 of 0.00 for virtually every method, with the sole
exception being XGB+SMOTE (0.20). This result indicates that nine training examples are
insufficient for any augmentation method to enable reliable minority-class recall, and that
the FOXA1 transcriptomic signature is not well-captured by interpolation from other subtypes.

In \textbf{COAD}, all classes achieve F1 $\geq$0.67 for most methods, and augmentation effects
are modest. CMS1 shows the most variation (0.67--0.84), primarily driven by confusion with CMS3.

\subsection{Confusion Matrices}

Figure~\ref{fig:cm} presents row-normalized confusion matrices for the six representative methods
(seed 42). In BRCA, the off-diagonal mass between Luminal A and Luminal B is consistent across
methods, reflecting biological proximity. cVAE introduces notable confusion into the HER2 class
(8\% misclassified as Basal) not present in other methods. In PRAD, all methods correctly
classify Other and ERG with high recall, but FOXA1 is misclassified as Other or SPOP regardless
of augmentation strategy.

\begin{figure}[htbp]
  \centering
  \includegraphics[width=0.95\linewidth]{cm_grid.png}
  \caption{Row-normalized confusion matrices for six augmentation methods on BRCA (top), COAD
    (middle), and PRAD (bottom), seed 42 test set. Darker diagonal indicates higher per-class
    recall.}
  \label{fig:cm}
\end{figure}
\FloatBarrier

\subsection{Effect of Augmentation Volume}

Figure~\ref{fig:sweep} shows BRCA accuracy and macro F1 as a function of the number of synthetic
samples added per minority class ($N_\text{synth} \in \{100, 500, 1000, 2000, 3000\}$). SMOTE's
macro F1 increases monotonically with $N_\text{synth}$, reaching 0.82 at $N_\text{synth}=3000$,
suggesting it benefits from aggressive oversampling. WGAN-GP and cVAE show degrading macro F1
as more synthetic samples are added, indicating that their generated samples are misaligned with
the true minority-class distribution. Mixup remains roughly flat, hovering near baseline.

\begin{figure}[htbp]
  \centering
  \includegraphics[width=0.85\linewidth]{sweep_brca.png}
  \caption{BRCA accuracy (left) and macro F1 (right) vs.\ number of synthetic samples per
    minority class ($N_\text{synth}$) for SMOTE, cVAE, WGAN-GP, and Mixup. The dashed line is
    the no-augmentation baseline. SMOTE improves monotonically; generative models degrade.}
  \label{fig:sweep}
\end{figure}
\FloatBarrier

\subsection{Synthetic Quality vs.\ Downstream Utility}

A central question in generative augmentation is whether better synthetic data, as measured by
distributional fidelity, leads to better classifier performance. Figure~\ref{fig:synth_quality}
plots downstream $\Delta$macro F1 against MMD (left) and gene-gene correlation preservation
(right) for all method/dataset combinations.

The scatter plots reveal no clear positive relationship between either quality metric and
downstream gain. Notably, BRCA/cVAE achieves low MMD (good distributional matching) but yields
the largest drop in macro F1 ($-$0.101). BRCA/BorderlineSMOTE, with high Pearson correlation
preservation ($\sim$0.55) and moderate MMD, produces the largest macro F1 gain ($+$0.033).
WGAN-GP achieves high correlation preservation on BRCA but negative downstream utility. These
results indicate that MMD and correlation preservation are insufficient proxies for augmentation
utility and should not be used as the sole evaluation criterion for synthetic data in this setting.

\begin{figure}[htbp]
  \centering
  \includegraphics[width=0.85\linewidth]{synth_quality_vs_delta_f1.png}
  \caption{Synthetic sample quality vs.\ downstream $\Delta$macro F1 over no-augmentation
    baseline. Left: distribution match (MMD, lower = more similar to real data). Right:
    gene-gene Pearson correlation preservation. Neither metric correlates with classifier benefit.}
  \label{fig:synth_quality}
\end{figure}
\FloatBarrier

% ============================================================
\section{Discussion}
% ============================================================

Our results paint a nuanced picture of when and why augmentation helps in cancer genomics.

\paragraph{Interpolation beats generation.} The consistent advantage of SMOTE-family methods over
cVAE and WGAN-GP is striking given the substantial additional complexity of generative models.
We hypothesize that this is partly a data-volume effect: with a few hundred minority-class
training examples, the cVAE and WGAN-GP are likely underfitting their generative distributions
and producing samples with spurious signal. SMOTE, by contrast, synthesizes within the convex
hull of observed minority-class examples, which may be a safer inductive bias for small cohorts.
BorderlineSMOTE's focus on decision boundary regions may specifically help the classifier learn
the LumA/LumB boundary in BRCA, which is the dominant source of misclassification.

\paragraph{Generative model degradation at scale.} The sweep experiment (Figure~\ref{fig:sweep})
shows that increasing $N_\text{synth}$ for cVAE and WGAN-GP \emph{degrades} performance, a
counterintuitive finding that suggests the generative models are adding noise rather than signal.
This is consistent with mode collapse or posterior collapse in the generative models, which can
produce low-diversity synthetic samples that over-represent certain gene expression patterns.

\paragraph{Dataset properties mediate augmentation benefit.} The benefit of augmentation is
largest on BRCA (severe imbalance, large $n$), absent on PRAD (moderate $n$, extreme FOXA1
minority), and marginal on COAD (smaller $n$ but balanced classes). This suggests that
augmentation is most useful when (a) the majority-to-minority ratio is large and (b) the
minority class has enough examples to enable meaningful interpolation or generation. FOXA1 at
$n{=}9$ appears to fall below this threshold for all methods tested.

\paragraph{Quality metrics are insufficient.} The lack of correlation between MMD or gene-gene
correlation and downstream $\Delta$F1 has direct practical implications: evaluating augmentation
methods solely by distributional fidelity is unreliable. Practitioners should always evaluate on
a downstream classification task with appropriate class-aware metrics.

\subsection{Limitations}

This study has several important limitations. First, we evaluate a single classifier architecture
(MLP) as the primary downstream model, and results may differ for other architectures such as
attention-based models or graph neural networks on biological networks. Second, all experiments
use bulk RNA-seq, and findings may not transfer to single-cell RNA-seq, proteomics, or multi-omic
settings. Third, the PRAD FOXA1 subtype ($n{=}9$) is likely too rare to draw reliable conclusions
from any augmentation experiment, and a larger collection effort would be needed to properly
evaluate this class. Fourth, our generative models are trained on the full gene expression space
($\sim$20{,}000 genes) without dimensionality reduction; it is possible that pre-training in a
reduced space (e.g., PCA or variational embedding) would improve their performance. Fifth, all
three cohorts are from TCGA and may not represent the full distribution of tumor transcriptomes
encountered in clinical practice.

% ============================================================
\section{Conclusion}
% ============================================================

We presented the first systematic benchmark of data augmentation strategies for cancer subtype
classification from bulk RNA-seq across three TCGA cohorts. Our main findings are: (1)
interpolation-based methods, particularly BorderlineSMOTE and ADASYN, outperform deep generative
approaches on imbalanced datasets; (2) generative models (cVAE, WGAN-GP) fail to improve and
frequently harm performance, with the harm growing as more synthetic samples are added; (3)
augmentation benefit is dataset-dependent, with BRCA benefiting most and PRAD not benefiting
from any method tested; and (4) synthetic sample quality as measured by MMD and gene-gene
correlation does not predict downstream classifier benefit. Taken together, these findings
recommend BorderlineSMOTE as a default augmentation strategy for bulk RNA-seq cancer subtype
classification, while cautioning against the assumption that more sophisticated generative
models yield proportionally better utility. Future work should investigate augmentation in
lower-dimensional learned feature spaces, explore alternative minority-class-aware objectives
for training generative models, and develop evaluation metrics that better predict downstream
utility from synthetic sample properties.

% ============================================================
\bibliographystyle{unsrt}
\begin{thebibliography}{99}

\bibitem{Parker2009}
Parker, J.S., et al. (2009).
Supervised Risk Predictor of Breast Cancer Based on Intrinsic Subtypes.
\textit{Journal of Clinical Oncology}, 27(8), 1160--1167.

\bibitem{Guinney2015}
Guinney, J., et al. (2015).
The consensus molecular subtypes of colorectal cancer.
\textit{Nature Medicine}, 21(11), 1350--1356.

\bibitem{Cancer2015}
Cancer Genome Atlas Research Network (2015).
The Molecular Taxonomy of Primary Prostate Cancer.
\textit{Cell}, 163(4), 1011--1025.

\bibitem{Shorten2019}
Shorten, C., \& Khoshgoftaar, T.M. (2019).
A survey on image data augmentation for deep learning.
\textit{Journal of Big Data}, 6(1), 1--48.

\bibitem{Wei2019}
Wei, J., \& Zou, K. (2019).
EDA: Easy Data Augmentation Techniques for Boosting Performance on Text Classification Tasks.
\textit{Proceedings of EMNLP-IJCNLP 2019}.

\bibitem{Chawla2002}
Chawla, N.V., et al. (2002).
SMOTE: Synthetic Minority Over-sampling Technique.
\textit{Journal of Artificial Intelligence Research}, 16, 321--357.

\bibitem{Han2005}
Han, H., Wang, W.Y., \& Mao, B.H. (2005).
Borderline-SMOTE: A New Over-Sampling Method in Imbalanced Data Sets Learning.
\textit{ICIC 2005}, Lecture Notes in Computer Science, vol. 3644.

\bibitem{He2008}
He, H., et al. (2008).
ADASYN: Adaptive synthetic sampling approach for imbalanced learning.
\textit{2008 IEEE International Joint Conference on Neural Networks (IJCNN)}.

\bibitem{Blagus2013}
Blagus, R., \& Lusa, L. (2013).
SMOTE for high-dimensional class-imbalanced data.
\textit{BMC Bioinformatics}, 14, 106.

\bibitem{Lemasson2022}
Lemasson, B., et al. (2022).
Oversampling strategies for imbalanced genomic datasets.
\textit{Briefings in Bioinformatics}, 23(4).

\bibitem{Kingma2014}
Kingma, D.P., \& Welling, M. (2014).
Auto-Encoding Variational Bayes.
\textit{Proceedings of ICLR 2014}.

\bibitem{Sohn2015}
Sohn, K., Lee, H., \& Yan, X. (2015).
Learning Structured Output Representation using Deep Conditional Generative Models.
\textit{NeurIPS 2015}.

\bibitem{Marouf2020}
Marouf, M., et al. (2020).
Realistic in silico generation and augmentation of single-cell RNA-seq data using generative
adversarial networks.
\textit{Nature Communications}, 11, 166.

\bibitem{Way2018}
Way, G.P., \& Greene, C.S. (2018).
Extracting a biologically relevant latent space from cancer transcriptomes with variational
autoencoders.
\textit{Pacific Symposium on Biocomputing}, 23, 80--91.

\bibitem{Goodfellow2014}
Goodfellow, I., et al. (2014).
Generative Adversarial Nets.
\textit{NeurIPS 2014}.

\bibitem{Gulrajani2017}
Gulrajani, I., et al. (2017).
Improved Training of Wasserstein GANs.
\textit{NeurIPS 2017}.

\bibitem{Zhang2018}
Zhang, H., et al. (2018).
Mixup: Beyond Empirical Risk Minimization.
\textit{Proceedings of ICLR 2018}.

\bibitem{Guo2019}
Guo, H., et al. (2019).
Mixup as locally linear out-of-manifold regularization.
\textit{Proceedings of AAAI 2019}.

\bibitem{Yoon2020}
Yoon, J., Jordon, J., \& van der Schaar, M. (2020).
GAIN: Missing Data Imputation using Generative Adversarial Nets.
\textit{ICML 2018}.

\bibitem{Levy2020}
Levy, J.J., et al. (2020).
Application of generative adversarial networks (GAN) for ophthalmology image domains.
\textit{Eye and Vision}, 7, 41.

\end{thebibliography}

\appendix

\section{Data Acquisition Deviations}
\label{app:data}

Three deviations from the original data acquisition specification are noted:

\textbf{1. Clinical files from GitHub LFS mirror.} The cBioPortal S3 datahub bucket
(\texttt{cbioportal-datahub.s3.amazonaws.com}) returned 403 Forbidden across all tested paths
as of 2026-04-27. Clinical files were retrieved from the public GitHub LFS mirror at
\texttt{media.githubusercontent.com/media/cBioPortal/datahub/master/public/}. Content is
identical to the S3 tarball.

\textbf{2. PRAD clinical from \texttt{prad\_tcga\_pub} (2015 Cell).} The Pan-Cancer Atlas PRAD
study (\texttt{prad\_tcga\_pan\_can\_atlas\_2018}) contains only the literal string ``PRAD'' in
its SUBTYPE column, not molecular subtypes. We use \texttt{prad\_tcga\_pub} (TCGA 2015 Cell
paper), which provides the expected subtype distribution. Pan-Cancer Atlas files are retained
for traceability.

\textbf{3. COAD CMS labels via public GitHub mirror.} The canonical source (Synapse syn4978511)
requires account registration. The same \texttt{cms\_labels\_public\_all.txt} file is available
via the Sage-Bionetworks \texttt{crc-cms-kras} GitHub repository and was used in its place.

\section{Hyperparameter Details}
\label{app:hyperparams}

\paragraph{MLP classifier.} Hidden dimensions: [1024, 512, 256]. Dropout: $p=0.3$. Batch
normalization after each hidden layer. Optimizer: Adam, learning rate $5\times10^{-4}$,
weight decay $10^{-5}$. Batch size: 64. Early stopping: patience 20 epochs on validation
macro F1. Maximum epochs: 200.

\paragraph{cVAE.} Encoder/decoder hidden dims: [512, 256, 128]. Latent dimension: 64.
KL weight annealing from 0 to 1 over 50 epochs. Optimizer: Adam, learning rate $10^{-3}$.
Training epochs: 100.

\paragraph{WGAN-GP.} Generator/critic hidden dims: [256, 512, 1024] / [1024, 512, 256].
Gradient penalty weight $\lambda=10$. Critic steps per generator step: 5. Optimizer: Adam,
$\alpha=10^{-4}$, $\beta_1=0$, $\beta_2=0.9$. Training epochs: 200.

\paragraph{SMOTE/variants.} $k=5$ nearest neighbors (default); tuned variant selects
$k \in \{3, 5, 7, 10\}$ by validation macro F1.

\paragraph{Mixup.} $\alpha=0.2$ (Beta distribution parameter). Applied during training.

\end{document}